\documentclass[11pt]{article}

\usepackage[T1]{fontenc}
\usepackage{lmodern}
\usepackage{amsmath,amssymb,amsthm,mathtools}
\usepackage{microtype}
\usepackage[hidelinks]{hyperref}
\usepackage{enumitem}
\usepackage[a4paper,margin=30mm]{geometry}
\usepackage{booktabs}

\newtheorem{theorem}{Theorem}[section]
\newtheorem{proposition}[theorem]{Proposition}
\newtheorem{lemma}[theorem]{Lemma}
\newtheorem{corollary}[theorem]{Corollary}
\theoremstyle{definition}
\newtheorem{definition}[theorem]{Definition}
\newtheorem{example}[theorem]{Example}
\theoremstyle{remark}
\newtheorem{remark}[theorem]{Remark}

\newcommand{\vp}{v_p}
\newcommand{\G}{G}
\newcommand{\Tp}{T_p}
\newcommand{\rp}{r_p}
\newcommand{\spdigit}{s_p}
\newcommand{\N}{\mathbb N}
\newcommand{\Cd}{C_d}

\title{Pascal Cofactors: Least Extremal Rows for Alternating Cofactors of Prime Powers}
\author{John Fairfax-Ball}
\date{}

\begin{document}
\maketitle

\begin{abstract}
For integers \(m\ge2\) and \(N>m\) with \(m\mid N\), define the restricted
binomial greatest common divisor
\[
  G(N;m)=\gcd\left\{\binom Nk:0<k<N,\ m\mid k\right\}.
\]
Fix a prime \(p\nmid m\), and let \(r_p(m)\) be the least positive \(r\) with
\(m<p^r\). Earlier work proves that \(r_p(m)\) is the global maximum of
\(v_p(G(N;m))\) over admissible rows and defines the least row \(T_p(m)\)
attaining it.

We study the alternating cofactors
\[
  C_d(p^a)=\frac{p^{ad}+1}{p^a+1},
  \qquad a\ge1,\quad d\ge3\text{ odd}.
\]
Writing \(Q=p^a\) and \(d=2s+1\), the central identity is the finite-window
scaling law
\[
  v_p\binom{C_d(Q)q}{C_d(Q)j}
  =as+v_p\binom qj
  \qquad(2\le q\le Q,\ 1\le j<q).
\]
It yields an exact pre-extremal restricted-GCD formula:
\[
  v_p(G(C_d(Q)q;C_d(Q)))
  =
  as+
  \begin{cases}
    1,&q\text{ is a positive power of }p,\\
    0,&\text{otherwise},
  \end{cases}
\]
for \(2\le q\le Q\).  A separate calculation at the sparse endpoint
\(q=Q+1\) gives valuation \(a(d-1)=r_p(C_d(Q))\).  Consequently, for every
odd \(d\ge5\),
\[
  T_p(C_d(p^a))=p^{ad}+1.
\]
For the cubic cofactor \(p^{2a}-p^a+1\), the exponent-one case behaves
differently:
\[
  T_p(p^{2a}-p^a+1)=
  \begin{cases}
    p(p^2-p+1),&a=1,\\
    p^{3a}+1,&a\ge2.
  \end{cases}
\]
The full theorem package has been formalised in Lean 4 and registered in the
Palomar registry.  The restricted-GCD family, Kummer's theorem, and the global
extremal framework are prior work; the paper records the relevant literature
boundaries explicitly.
\end{abstract}

\medskip
\noindent\textbf{2020 Mathematics Subject Classification.} 11B65.

\noindent\textbf{Keywords.} binomial coefficients; greatest common divisor;
\(p\)-adic valuation; Kummer's theorem; Pascal triangle; Lean.

\section{Introduction}\label{sec:intro}

Divisibility in Pascal's triangle is governed locally by the base-\(p\)
arithmetic of the upper and lower indices.  Kummer's theorem identifies
\(v_p\binom nk\) with a carry count in base \(p\), while Legendre's formula
gives the equivalent digit-sum expression
\[
  v_p\binom nk
  =
  \frac{s_p(k)+s_p(n-k)-s_p(n)}{p-1}.
\]
These classical descriptions are especially effective when the selected lower
indices have a rigid arithmetic shape.

We consider a fixed divisibility restriction.  For \(m\ge2\) and \(N>m\)
with \(m\mid N\), let
\[
  G(N;m)=\gcd\left\{\binom Nk:0<k<N,\ m\mid k\right\}.
\]
Writing \(N=mq\), this is the same fixed-multiple family denoted \(g(m,q)\)
by Wu \cite{Wu2026}; the family itself is therefore not new.  McTague studied
the same selector under congruence hypotheses on the valuation prime
\cite{McTague2017}.  Other work studies related but different restrictions,
including coprime, non-coprime, and \(\gcd(N,k)=d\) selectors
\cite{Hong2016,ChiuYuanZhou2023,ChungYangZhou2025,ChungYang2026}.

The row-extremal framework used here comes from the author's preceding work
\cite{FairfaxBallExtremes2026}.  For a prime \(p\nmid m\), put
\[
  r_p(m)=\min\{r\ge1:m<p^r\}.
\]
The predecessor theorem states that the largest possible value of
\(v_p(G(N;m))\), over all admissible rows \(N\), is exactly \(r_p(m)\), and
that an extremal row exists.  Only after this existence result is in place does
one define the least extremal row
\[
  T_p(m)=
  \min\{N>m:m\mid N,\ v_p(G(N;m))=r_p(m)\}.
\]

This paper studies \(T_p(m)\) when \(m\) is the alternating cofactor in
\[
  Q^d+1=(Q+1)C_d(Q),
  \qquad
  C_d(Q)=\frac{Q^d+1}{Q+1},
\]
where \(Q=p^a\) and \(d\ge3\) is odd.  Thus
\[
  C_d(Q)=Q^{d-1}-Q^{d-2}+\cdots-Q+1.
\]
For \(d=3\), the two factors are
\[
  Q+1,\qquad Q^2-Q+1=\Phi_6(Q),
\]
and the predecessor paper determines the least extremal row for the first
factor when \(a\ge2\):
\[
  T_p(Q+1)=Q^3+1.
\]
The cubic result below gives the complementary statement for the second factor,
but it is not a formal consequence of the factorisation.  Minimality for one
restriction modulus does not transfer to the complementary factor.

The proof is driven by a finite-window digit identity.  If \(d=2s+1\) and
\(m=C_d(Q)\), then for every \(1\le t\le Q\),
\[
  s_p(mt)=s_p(t)+as(p-1).
\]
Substituting this in Legendre's digit-sum formula gives the exact coefficient
scaling
\[
  v_p\binom{mq}{mj}
  =
  as+v_p\binom qj
  \qquad(2\le q\le Q,\ 1\le j<q).
\]
Minimising over \(j\) gives a complete row-by-row formula for the restricted
GCD before the sparse endpoint.  The endpoint \(q=Q+1\) lies just outside the
scaling window and is handled separately using the special row
\(Q^d+1=p^{ad}+1\).

The resulting picture has a useful dichotomy.  For every odd \(d\ge5\), all
multipliers \(2\le q\le Q\) remain strictly below the global extremal threshold
and \(q=Q+1\) is the first extremal multiplier.  In the cubic case \(d=3\),
the same conclusion holds for \(a\ge2\), but when \(a=1\) the multiplier
\(q=p\) already reaches the threshold.  This produces the complete cubic
classification stated in the abstract.

The paper is organised as follows.  Section~\ref{sec:setup} records the
restricted-GCD and extremal framework.  Section~\ref{sec:threshold} determines
the exact threshold for the cofactor family.  Sections~\ref{sec:digits} and
\ref{sec:scaling} prove the digit-sum and coefficient-scaling laws.
Section~\ref{sec:gcdwindow} derives the exact pre-extremal GCD formula.
Section~\ref{sec:endpoint} treats the sparse endpoint.
Section~\ref{sec:least} proves the least-row theorems, and
Section~\ref{sec:examples} gives illustrative cases.  The literature and
formal-verification boundaries are recorded in
Sections~\ref{sec:related} and \ref{sec:formal}.

\section{Restricted GCDs and the extremal framework}\label{sec:setup}

Throughout, \(p\) is prime.  For a positive integer \(n\), \(v_p(n)\) denotes
the exponent of \(p\) in \(n\), and \(s_p(n)\) denotes the sum of the base-\(p\)
digits of \(n\).

\begin{definition}
For integers \(m\ge2\) and \(N>m\) with \(m\mid N\), define
\[
  G(N;m)=\gcd\left\{\binom Nk:0<k<N,\ m\mid k\right\}.
\]
For prime \(p\nmid m\), define
\[
  r_p(m)=\min\{r\ge1:m<p^r\}.
\]
\end{definition}

If \(N=mq\) with \(q\ge2\), the selected lower indices are precisely
\(m,2m,\ldots,(q-1)m\).  Therefore
\begin{equation}\label{eq:gcd-min}
  v_p(G(mq;m))
  =
  \min_{1\le j<q}v_p\binom{mq}{mj}.
\end{equation}

We use Kummer's theorem and its digit-sum reformulation
\cite{Kummer1852,Granville1997}:
\begin{equation}\label{eq:kummer-digit}
  v_p\binom nk
  =
  \frac{s_p(k)+s_p(n-k)-s_p(n)}{p-1}.
\end{equation}

The following theorem is prior work and supplies the global meaning of the
threshold \(r_p(m)\).

\begin{theorem}[Global extremal theorem, \cite{FairfaxBallExtremes2026}]
\label{thm:prior-global}
Let \(p\) be prime and \(m\ge2\) with \(p\nmid m\).  Then
\[
  \max_{\substack{N>m\\m\mid N}}v_p(G(N;m))=r_p(m).
\]
Moreover the extremal set is nonempty, so the least extremal row
\[
  T_p(m)=
  \min\{N>m:m\mid N,\ v_p(G(N;m))=r_p(m)\}
\]
is well defined.
\end{theorem}

Our task is therefore not to reprove the global maximum, but to identify the
first row attaining it for the cofactor family.

\section{Alternating cofactors and the exact threshold}\label{sec:threshold}

Fix \(a\ge1\), put \(Q=p^a\), and let \(d=2s+1\ge3\) be odd.  Set
\[
  m=C_d(Q)=\frac{Q^d+1}{Q+1}.
\]
The alternating expansion is
\[
  m=Q^{2s}-Q^{2s-1}+\cdots-Q+1.
\]

\begin{lemma}[Positive block expansion]\label{lem:block}
One has
\[
  m=1+(Q-1)\sum_{i=0}^{s-1}Q^{2i+1}.
\]
\end{lemma}

\begin{proof}
For each \(i\),
\[
  (Q-1)Q^{2i+1}=Q^{2i+2}-Q^{2i+1}.
\]
Summing these pairs and adding \(1\) gives the alternating expansion of
\(m\).
\end{proof}

\begin{lemma}[Prime nondivisibility]\label{lem:coprime}
One has \(m\equiv1\pmod p\).  In particular \(p\nmid m\).
\end{lemma}

\begin{proof}
Since \(Q=p^a\equiv0\pmod p\), every nonconstant term of the alternating
expansion vanishes modulo \(p\).
\end{proof}

\begin{proposition}[Exact threshold interval]\label{prop:threshold}
For every odd \(d\ge3\),
\[
  p^{a(d-1)-1}<C_d(p^a)<p^{a(d-1)}.
\]
Consequently,
\[
  r_p(C_d(p^a))=a(d-1)=2as.
\]
\end{proposition}

\begin{proof}
By Lemma~\ref{lem:block},
\[
  m>(Q-1)Q^{d-2}.
\]
Since \(Q=p^a\) and
\[
  p^a-1\ge p^{a-1},
\]
we obtain
\[
  m>p^{a-1}p^{a(d-2)}
   =p^{a(d-1)-1}.
\]

For the upper bound, use
\[
  (Q+1)m=Q^d+1.
\]
Because \(Q^{d-1}>1\),
\[
  Q^d+1<Q^d+Q^{d-1}=(Q+1)Q^{d-1},
\]
hence \(m<Q^{d-1}=p^{a(d-1)}\).

Thus \(a(d-1)\) is the least positive exponent \(r\) with \(m<p^r\).
\end{proof}

\begin{remark}
The strict lower bound in Proposition~\ref{prop:threshold} is one power of
\(p\) below the upper threshold.  This exact interval is important: the
opposite lower inequality \(p^{a(d-1)}<m\) is false.
\end{remark}

\section{The finite-window digit-sum identity}\label{sec:digits}

The key structural fact is that multiplication by \(m=C_d(Q)\) adds a constant
amount to the base-\(p\) digit sum throughout the complete interval
\(1\le t\le Q\).

\begin{lemma}[Block complement]\label{lem:complement}
If \(0\le u<Q=p^a\), then
\[
  s_p((Q-1)-u)=a(p-1)-s_p(u).
\]
\end{lemma}

\begin{proof}
The number \(Q-1=p^a-1\) has \(a\) base-\(p\) digits, all equal to \(p-1\).
Subtracting \(u<Q\) from \(Q-1\) is therefore digitwise complementation with no
borrow.  Each digit \(u_i\) is replaced by \(p-1-u_i\).
\end{proof}

\begin{theorem}[Exact digit-sum shift]\label{thm:digit-shift}
For every \(1\le t\le Q\),
\[
  s_p(mt)=s_p(t)+as(p-1).
\]
\end{theorem}

\begin{proof}
Assume first that \(1\le t<Q\).  Starting from
Lemma~\ref{lem:block}, one may normalise \(mt\) in base \(Q\) as
\begin{equation}\label{eq:baseQ}
  mt
  =
  t
  +\sum_{i=0}^{s-1}(Q-t)Q^{2i+1}
  +\sum_{i=1}^{s}(t-1)Q^{2i}.
\end{equation}
Indeed,
\[
  (Q-t)Q^{2i+1}+(t-1)Q^{2i+2}
  =
  t(Q-1)Q^{2i+1}.
\]
All displayed base-\(Q\) digits lie in \(\{0,\ldots,Q-1\}\).  Since
\(Q=p^a\), distinct base-\(Q\) positions occupy disjoint blocks of \(a\)
base-\(p\) digits.  Hence
\[
  s_p(mt)
  =
  s_p(t)+s\bigl(s_p(Q-t)+s_p(t-1)\bigr).
\]
Now \(Q-t=(Q-1)-(t-1)\), so Lemma~\ref{lem:complement} gives
\[
  s_p(Q-t)+s_p(t-1)=a(p-1),
\]
and the result follows.

For \(t=Q\), multiply Lemma~\ref{lem:block} by \(Q\):
\[
  Qm
  =
  Q+(Q-1)(Q^2+Q^4+\cdots+Q^{2s}).
\]
Its base-\(p\) digit sum is \(1+sa(p-1)\), while \(s_p(Q)=1\).
\end{proof}

The endpoint \(t=Q\) is included in the same identity, but the later binomial
scaling theorem only uses triples \(j,q-j,q\) lying in \([1,Q]\); this is why
the multiplier range terminates at \(q=Q\).

\section{Coefficient scaling in the pre-endpoint window}\label{sec:scaling}

\begin{theorem}[Finite-window coefficient scaling]\label{thm:scaling}
Let \(2\le q\le Q\) and \(1\le j<q\).  Then
\[
  v_p\binom{mq}{mj}
  =
  as+v_p\binom qj.
\]
\end{theorem}

\begin{proof}
The integers \(j\), \(q-j\), and \(q\) all lie in \([1,Q]\).  Apply
Theorem~\ref{thm:digit-shift} to all three and substitute into
\eqref{eq:kummer-digit}:
\[
\begin{aligned}
  v_p\binom{mq}{mj}
  &=
  \frac{s_p(mj)+s_p(m(q-j))-s_p(mq)}{p-1}\\
  &=
  \frac{
    s_p(j)+as(p-1)
    +s_p(q-j)+as(p-1)
    -s_p(q)-as(p-1)}
  {p-1}\\
  &=
  as+v_p\binom qj.
\end{aligned}
\]
\end{proof}

Taking \(j=1\) gives an exact selected-coefficient formula.

\begin{corollary}[Selected coefficient]\label{cor:selected}
For \(2\le q\le Q\),
\[
  v_p\binom{mq}{m}=as+v_p(q).
\]
Hence
\[
  \max_{2\le q\le Q}v_p\binom{mq}{m}
  =
  as+a
  =
  \frac{a(d+1)}2,
\]
and the maximum occurs uniquely at \(q=Q=p^a\).
\end{corollary}

\begin{proof}
Since \(\binom q1=q\), the first formula is
Theorem~\ref{thm:scaling} with \(j=1\).  For \(q\le p^a\),
\(v_p(q)\le a\), with equality exactly when \(q=p^a\).
\end{proof}

\begin{remark}\label{rem:selected-not-gcd}
A single selected coefficient need not control extremality of the restricted
GCD.  In the cubic case \(d=3\), the maximum in
Corollary~\ref{cor:selected} equals \(2a=r_p(m)\), even though for \(a\ge2\)
no row with \(q\le Q\) is extremal.  The distinction is resolved by minimising
over all selected coefficients.
\end{remark}

\section{The exact pre-extremal restricted-GCD formula}
\label{sec:gcdwindow}

For \(q\ge2\), write
\[
  \mu_p(q)=\min_{1\le j<q}v_p\binom qj.
\]

\begin{lemma}[Interior minimum]\label{lem:mu}
For every \(q\ge2\),
\[
  \mu_p(q)=
  \begin{cases}
    1,&q=p^b\text{ for some }b\ge1,\\
    0,&q\text{ is not a power of }p.
  \end{cases}
\]
\end{lemma}

\begin{proof}
Suppose first that \(q=p^b\).  Every \(1\le j<q\) forces at least one borrow
when \(j\) is subtracted from the base-\(p\) expansion \(10\cdots0\), so
Kummer gives \(v_p\binom qj\ge1\).  Taking \(j=p^{b-1}\), the addition
\[
  p^{b-1}+(p-1)p^{b-1}=p^b
\]
has exactly one carry.  Thus the minimum is \(1\).

Now suppose \(q\) is not a power of \(p\), and let \(b=v_p(q)\).  The base-\(p\)
digit of \(q\) at position \(b\) is nonzero, and \(p^b<q\).  Taking
\(j=p^b\), the subtraction \(q-j\) requires no borrow.  Hence
\(v_p\binom q{p^b}=0\).
\end{proof}

\begin{theorem}[Exact pre-endpoint GCD valuation]\label{thm:gcd-window}
For \(2\le q\le Q\),
\[
  v_p(G(mq;m))
  =
  as+
  \begin{cases}
    1,&q=p^b\text{ for some }1\le b\le a,\\
    0,&\text{otherwise}.
  \end{cases}
\]
\end{theorem}

\begin{proof}
By \eqref{eq:gcd-min} and Theorem~\ref{thm:scaling},
\[
\begin{aligned}
  v_p(G(mq;m))
  &=
  \min_{1\le j<q}v_p\binom{mq}{mj}\\
  &=
  as+\min_{1\le j<q}v_p\binom qj\\
  &=
  as+\mu_p(q).
\end{aligned}
\]
Apply Lemma~\ref{lem:mu}.  Since \(q\le Q=p^a\), the positive powers of \(p\)
in the interval are exactly \(p,p^2,\ldots,p^a\).
\end{proof}

\begin{corollary}[Pre-endpoint maximum and argmax]\label{cor:gcd-max}
For every odd \(d\ge3\),
\[
  \max_{2\le q\le Q}v_p(G(mq;m))
  =
  as+1
  =
  \frac{a(d-1)}2+1.
\]
The complete argmax set is
\[
  \{p,p^2,\ldots,p^a\}.
\]
\end{corollary}

The power/non-power dichotomy in Lemma~\ref{lem:mu} is classical for the
ordinary binomial row.  What is special here is that
Theorem~\ref{thm:scaling} transports it, with the additive shift \(as\), to
the cofactor-restricted rows.

\section{The sparse endpoint}\label{sec:endpoint}

The multiplier \(q=Q+1\) is outside the finite scaling window.  It is also
arithmetically distinguished:
\[
  m(Q+1)=Q^d+1=p^{ad}+1.
\]

\begin{lemma}[A \(p^L+1\) valuation identity]\label{lem:plus-one}
Let \(P=p^L\) with \(L\ge1\), and let \(1\le k<P\).  Then
\[
  v_p\binom{P+1}{k}
  =
  L-v_p(k)-v_p(P+1-k).
\]
\end{lemma}

\begin{proof}
The identity
\[
  k\binom Pk=P\binom{P-1}{k-1}
\]
and Kummer's theorem give
\[
  v_p\binom Pk=L-v_p(k),
\]
because \(P-1\) has all lower \(L\) base-\(p\) digits equal to \(p-1\), so
\(\binom{P-1}{k-1}\) has valuation zero.  Next,
\[
  (P+1-k)\binom{P+1}{k}=(P+1)\binom Pk.
\]
Since \(p\nmid P+1\), taking \(p\)-adic valuations yields the formula.
\end{proof}

\begin{theorem}[Endpoint coefficients and endpoint GCD]
\label{thm:endpoint}
For \(1\le j\le Q\),
\[
  v_p\binom{m(Q+1)}{mj}
  =
  ad-v_p(j)-v_p(Q+1-j).
\]
Consequently,
\[
  v_p(G(m(Q+1);m))
  =
  ad-a
  =
  a(d-1)
  =
  r_p(m).
\]
\end{theorem}

\begin{proof}
Put \(P=Q^d=p^{ad}\).  By Proposition~\ref{prop:threshold},
\[
  1\le mj\le mQ<Q^d=P
\]
for \(1\le j\le Q\).  Lemma~\ref{lem:plus-one} applies with \(k=mj\).
Since \(p\nmid m\),
\[
  v_p(mj)=v_p(j).
\]
Also
\[
  P+1-mj
  =
  m(Q+1)-mj
  =
  m(Q+1-j),
\]
hence
\[
  v_p(P+1-mj)=v_p(Q+1-j).
\]
This proves the coefficient formula.

Now
\[
  j+(Q+1-j)=Q+1\equiv1\pmod p,
\]
so \(j\) and \(Q+1-j\) cannot both be divisible by \(p\).  Since each lies in
\([1,Q]\),
\[
  v_p(j)+v_p(Q+1-j)\le a.
\]
Therefore every selected endpoint coefficient has valuation at least
\(ad-a\).  Equality occurs at \(j=1\), for which \(Q+1-j=Q\), and also at
\(j=Q\).  Taking the minimum over the selected coefficients gives
\[
  v_p(G(m(Q+1);m))=ad-a=a(d-1).
\]
The final equality with \(r_p(m)\) is Proposition~\ref{prop:threshold}.
\end{proof}

Thus the sparse row \(Q^d+1\) is always extremal.  The remaining question is
whether it is the first extremal row.

\section{Least extremal rows}\label{sec:least}

\subsection{Odd degree \(d\ge5\)}

\begin{theorem}[Least extremal row for alternating cofactors]
\label{thm:mainA}
Let \(p\) be prime, \(a\ge1\), and \(d\ge5\) odd.  Then
\[
  T_p(C_d(p^a))=p^{ad}+1.
\]
Equivalently, the least extremal multiplier is \(q=p^a+1\).
\end{theorem}

\begin{proof}
Write \(d=2s+1\).  Since \(d\ge5\), we have \(s\ge2\).  By
Proposition~\ref{prop:threshold},
\[
  r_p(m)=2as.
\]
For every \(2\le q\le Q\), Theorem~\ref{thm:gcd-window} gives
\[
  v_p(G(mq;m))\le as+1.
\]
Because \(a\ge1\) and \(s\ge2\),
\[
  as+1<2as.
\]
Hence no multiplier \(2\le q\le Q\) is extremal.  On the other hand,
Theorem~\ref{thm:endpoint} shows that \(q=Q+1\) is extremal.  Therefore it is
the least extremal multiplier, and
\[
  T_p(m)=m(Q+1)=Q^d+1=p^{ad}+1.
\]
\end{proof}

\subsection{The cubic cofactor}

When \(d=3\), one has \(s=1\) and
\[
  m=C_3(Q)=Q^2-Q+1=p^{2a}-p^a+1,
  \qquad
  r_p(m)=2a.
\]
The inequality used above becomes \(a+1<2a\), so the exponent-one case must be
separated.

\begin{theorem}[Complete cubic classification]\label{thm:mainC}
For every prime \(p\) and every \(a\ge1\),
\[
  T_p(p^{2a}-p^a+1)
  =
  \begin{cases}
    p(p^2-p+1),&a=1,\\
    p^{3a}+1,&a\ge2.
  \end{cases}
\]
\end{theorem}

\begin{proof}
Assume first that \(a=1\).  Then \(Q=p\), \(m=p^2-p+1\), and
\(r_p(m)=2\).  Theorem~\ref{thm:gcd-window} gives
\[
  v_p(G(mq;m))
  =
  1+
  \begin{cases}
    1,&q\text{ is a positive power of }p,\\
    0,&\text{otherwise}
  \end{cases}
  \qquad(2\le q\le p).
\]
The only positive power of \(p\) in this interval is \(q=p\).  Thus every
\(2\le q<p\) has valuation \(1<2\), whereas
\[
  v_p(G(mp;m))=2=r_p(m).
\]
Hence the least extremal multiplier is \(p\), and
\[
  T_p(m)=mp=p(p^2-p+1).
\]

Now let \(a\ge2\).  For \(2\le q\le Q\), Theorem~\ref{thm:gcd-window} gives
\[
  v_p(G(mq;m))\le a+1<2a=r_p(m).
\]
Therefore no pre-endpoint multiplier is extremal.  The endpoint
Theorem~\ref{thm:endpoint} gives equality at \(q=Q+1\), so
\[
  T_p(m)=m(Q+1)=Q^3+1=p^{3a}+1.
\]
\end{proof}

\begin{remark}[Complementary-factor predecessor]\label{rem:complementary}
For \(a\ge2\), the predecessor theorem \cite{FairfaxBallExtremes2026} states
\[
  T_p(p^a+1)=p^{3a}+1.
\]
Since
\[
  (p^a+1)(p^{2a}-p^a+1)=p^{3a}+1,
\]
Theorem~\ref{thm:mainC} shows that the same sparse row is least extremal for
both complementary factors when \(a\ge2\).  The two least-row results are
nevertheless logically distinct: the predecessor theorem does not imply the
cofactor-side lower-row non-attainment proved here.
\end{remark}

\begin{remark}[The edge \(p=2,a=1\)]
For \(p=2\) and \(a=1\), Theorem~\ref{thm:mainC} gives \(m=3\) and
\(T_2(3)=6\).  The valuation \(v_2(G(6;3))=2\) is a known isolated instance,
explicitly recorded by McTague \cite{McTague2017} and by the author's earlier
minus-one work \cite{FairfaxBallMinusOne2026}.  The present theorem includes
that case as part of the family classification; it is not presented as a new
isolated observation.
\end{remark}

\section{Examples and structural features}\label{sec:examples}

\begin{example}[A quintic cofactor]
Take \(p=2\), \(a=1\), \(d=5\).  Then \(Q=2\),
\[
  m=C_5(2)=\frac{2^5+1}{2+1}=11,
  \qquad
  r_2(11)=4.
\]
The only pre-endpoint multiplier is \(q=2\), and
Theorem~\ref{thm:gcd-window} gives
\[
  v_2(G(22;11))=as+1=1\cdot2+1=3.
\]
At \(q=3\),
\[
  m(Q+1)=33=2^5+1,
\]
and Theorem~\ref{thm:endpoint} gives valuation \(4\).  Hence
\[
  T_2(11)=33.
\]
\end{example}

\begin{example}[Why the selected coefficient is insufficient]
Take \(p=3\), \(a=2\), \(d=3\).  Then \(Q=9\),
\[
  m=9^2-9+1=73,
  \qquad
  r_3(73)=4.
\]
At the final pre-endpoint multiplier \(q=9\),
Corollary~\ref{cor:selected} gives
\[
  v_3\binom{657}{73}=a+v_3(9)=4,
\]
already equal to the global threshold.  Nevertheless
Theorem~\ref{thm:gcd-window} gives
\[
  v_3(G(657;73))=a+1=3,
\]
because another selected coefficient has smaller valuation.  The first
extremal restricted GCD occurs at
\[
  q=10,\qquad N=730=3^6+1.
\]
\end{example}

\begin{example}[The exponent-one cubic branch]
Take \(p=3\), \(a=1\).  Then
\[
  m=3^2-3+1=7,\qquad r_3(7)=2.
\]
The multiplier \(q=3\) already gives
\[
  v_3(G(21;7))=2,
\]
so \(T_3(7)=21\).  The sparse endpoint \(q=Q+1=4\) gives the later extremal
row \(28=3^3+1\), showing exactly why \(a=1\) must be separated from
\(a\ge2\) in Theorem~\ref{thm:mainC}.
\end{example}

The proof also identifies the full pre-endpoint extremal pattern within the
window: among \(2\le q\le Q\), the restricted-GCD valuation is higher by one
exactly at powers of \(p\).  By contrast, the selected coefficient
\(\binom{mq}{m}\) increases according to \(v_p(q)\) and has a unique maximum
at \(q=Q\).  These two statistics agree on neither their maxima nor their
argmax sets in general.

\section{Relation to prior work}\label{sec:related}

The fixed-multiple restricted-GCD family has a substantial prior literature,
and the theorem package above should be read within that context.

McTague studies
\[
  \gcd_{0<k<n/q}\binom n{qk},
\]
which becomes \(G(N;m)\) after renaming \(n=N\) and \(q=m\).  His theorem under
\(p\equiv1\pmod m\), together with a corrected same-residue extension, is not
applicable to the present family in general because the cofactor modulus is
typically larger than the valuation prime.  His paper does contain the
specific valuation \(v_2(G(6;3))=2\) \cite{McTague2017}.

Wu uses exactly the notation-equivalent family
\[
  g(m,n)=\gcd\left\{\binom{mn}{mk}:1\le k<n\right\}.
\]
For a restriction modulus that is itself a power of the valuation prime, Wu
obtains the familiar power/non-power dichotomy in \(n\) \cite{Wu2026}.  In the
present cofactor family, however, \(p\nmid m\) by Lemma~\ref{lem:coprime}; the
hypotheses are therefore different.  The ordinary-row minimum
Lemma~\ref{lem:mu} has the same dichotomy, but
Theorem~\ref{thm:scaling} is what transfers it to cofactor-restricted rows.

Other papers select lower indices by different arithmetic conditions.  Hong
uses a coprimality selector \cite{Hong2016}.  Chiu--Yuan--Zhou study
non-coprime indices in a central band \cite{ChiuYuanZhou2023}, while
Chung--Yang--Zhou use selectors of the form \(\gcd(N,k)=d\)
\cite{ChungYangZhou2025}.  Guo--Qiu--Cao--Feng--Gao take a GCD across the row
multiplier for a fixed lower index rather than a per-row GCD over lower-index
multiples \cite{GuoEtAl2026}.  These are adjacent aggregation problems rather
than notation-equivalent versions of the present least-row question.

The author's preceding papers are also genuine prior work.  The minus-one
paper includes a common-\(p^c\) scaling theorem
\cite{FairfaxBallMinusOne2026}; it removes a common power of \(p\) from the row
and restriction modulus and is different from the additive scaling in
Theorem~\ref{thm:scaling}.  The extremal paper
\cite{FairfaxBallExtremes2026} supplies Theorem~\ref{thm:prior-global}, the
least-row framework, and the complementary theorem
\(T_p(p^a+1)=p^{3a}+1\) for \(a\ge2\).

A theorem-level literature audit was carried out against the final statements
proved here, with a cutoff of 2 October 2026.  No mathematically equivalent
or theorem-level source implying the family-level coefficient-scaling,
pre-endpoint GCD, sparse-endpoint, or least-row statements was located in that
documented search.  This negative-search result is not a claim of historical
priority.  In particular, the recently published paper of Chung and Yang
\cite{ChungYang2026} remained subscription-inaccessible at theorem level
during the audit.  Its abstract concerns non-coprime indices, a different
visible selector, but inaccessible internal results could not be ruled out as
overlapping special cases.  That comparison therefore remains explicitly
unresolved.

\section{Formal verification}\label{sec:formal}

The complete theorem package was formalised in Lean 4 before this manuscript
was written.  The formal development uses Lean
\texttt{v4.35.0-rc2} and a pinned Mathlib dependency through the author's
Pascal Extremes formalisation.  The latter supplies prior infrastructure for
\(G\), \(r_p\), extremal rows, and \(T_p\); the cofactor-specific threshold,
digit arithmetic, scaling, exact GCD formula, endpoint theorem, and least-row
results are proved in the present repository.

The principal mathematical statements correspond to the following public Lean
declarations:
\begin{center}
\begin{tabular}{@{}ll@{}}
\toprule
Mathematical statement & Lean declaration\\
\midrule
Finite-window coefficient scaling &
\texttt{PascalCofactors.coefficient\_scaling}\\
Exact bounded GCD valuation &
\texttt{PascalCofactors.restricted\_gcd\_valuation\_bounded}\\
Odd \(d\ge5\) least-row theorem &
\texttt{PascalCofactors.targetA\_d}\\
Cubic classification &
\texttt{PascalCofactors.targetC}\\
\bottomrule
\end{tabular}
\end{center}

The formal sources contain no substantive \texttt{sorry} or \texttt{admit}.
The advertised Palomar Challenge was independently checked by Comparator and
the con-ron, NanoDa, and Lean kernels.  The package was registered as
\href{https://palomar-registry.org/entry.html?id=PALOMAR-2026-10-02-000014&version=1}
{\texttt{PALOMAR-2026-10-02-000014}, version 1},
with registered source commit
\texttt{02a52e71a0a5ab1e77824490d0c47650d4a45691}.

Formal verification and registry preservation establish that the stated Lean
theorems check under the recorded environment.  They do not establish
historical novelty, literature completeness, or independent human peer review.

\section{Concluding remarks}

For alternating cofactors
\[
  C_d(p^a)=\frac{p^{ad}+1}{p^a+1},
\]
the base-\(p\) structure of the cofactor gives an unusually rigid finite
window.  Multiplication by the cofactor shifts digit sums by the constant
\(a(d-1)(p-1)/2\), and this converts directly into an additive scaling law for
all selected binomial coefficients.  Minimising that identity produces the
exact row-by-row restricted-GCD valuation up to \(q=p^a\).

The next multiplier, \(q=p^a+1\), is governed by a different mechanism: the
row collapses to the sparse number \(p^{ad}+1\).  Comparing the pre-endpoint
maximum with the global threshold then identifies the least extremal row for
all odd \(d\ge5\), while the cubic case exposes a single structural
exceptional regime at exponent \(a=1\).

The complementary cubic factors
\[
  p^a+1
  \quad\text{and}\quad
  p^{2a}-p^a+1
\]
therefore share the least extremal row \(p^{3a}+1\) when \(a\ge2\), but for
different restriction moduli and with separate minimality arguments.  This
cofactor viewpoint suggests that other sparse factorizations of
\(p^L\pm1\) may support similarly rigid digit windows, though no general
classification is asserted here.

\paragraph{Acknowledgements.}
The author thanks Carl McTague and Chai Wah Wu for prior work on restricted
binomial GCDs, and Ernst Eduard Kummer for the classical valuation theorem on
which the carry and digit-sum arguments rest.  The Lean formalisation uses
Mathlib and builds on the author's Pascal Extremes formal infrastructure.
ChatGPT (OpenAI, GPT-5.6 Sol) assisted with proof development, literature
auditing, Lean formalisation, repository operations, and manuscript drafting
under human direction and review.

\bibliographystyle{plain}
\bibliography{references}

\end{document}
